\begin{frame}[standout]

    Focus sur MongoDB

    \note{
        MongoDB est l'exemple de référence pour illustrer les bases document.
        Le modèle BSON reprend la logique JSON tout en ajoutant des types adaptés
        au stockage et aux index.
    }

\end{frame}

\begin{frame}{Étude de cas}

    https://www.mongodb.com/solutions/customer-case-studies

    https://www.mongodb.com/solutions/customer-case-studies/sncfconnect

    https://www.mongodb.com/solutions/customer-case-studies/decathlon

    https://www.mongodb.com/solutions/customer-case-studies/anfsi

    https://www.mongodb.com/solutions/customer-case-studies/carrefour

    https://www.mongodb.com/solutions/customer-case-studies/loreal

    Uber et MongoDB (gestion des trajets en temps réel).

    Problématique.
    Solution NoSQL choisie.
    Architecture mise en place.
    Résultats (performances, coûts).

    \note{
        Pour chaque entreprise, faire le lien entre le besoin et les propriétés
        de MongoDB : documents proches du modèle métier, schéma évolutif,
        indexation et réplication. Demander quelle partie des données serait
        regroupée dans un document et laquelle nécessiterait une référence.
        Les liens servent de supports de préparation, pas de contenu à lire
        intégralement pendant le cours.
    }

\end{frame}
